% fixtures/uncited-assertions-sample.tex
% Contract fixture for the uncited-assertion checker — documents exactly what MUST and MUST NOT flag.
% Since 2026-08-27 that checker is the ESLint rule `paper/uncited-assertion` (eslint-rules/paper-claims.mjs);
% this file is its test, run by eslint-rules/paper-claims.harness.mjs.
% Each numbered case is annotated with the expected verdict. Keep in the repo (it IS the contract).
\documentclass{article}
\begin{document}

% (1) MUST FLAG — uncited empirical claim about prior work (percentage + verb, no cite, no "we").
Transformers outperform RNNs by 12\%.

% (2) MUST NOT FLAG — same claim, now carrying a citation.
Transformers outperform RNNs by 12\% \cite{vaswani2017}.

% (3) MUST NOT FLAG — float pointer; the verb refers to a Table, not to prior work.
Table 2 shows the aggregated results across all runs.

% (4) MUST NOT FLAG — the only number is a year (year guard).
In 2026, the field continued to attract new benchmarks.

% (5) MUST NOT FLAG — the paper's OWN reported result (we/our guard); numbers allowed uncited.
We find a 6\% reduction in false positives (p<0.05) over the baseline.

% (6) MUST NOT FLAG — a definition sentence, even though it contains a number.
A transformer block is defined as a stack of 6 attention sub-layers.

% (7) MUST NOT FLAG — version-number guard (v2.1) and a math-mode number ($n = 512$) both skipped;
%     no "we", so this isolates the version + math guards (no empirical signal survives).
The baseline model ships as release v2.1 with a hidden size of $n = 512$.

% (8) MUST FLAG — bare fuzzy quantifier making a claim about prior work, no cite.
Most prior detectors silently miss adversarial inputs.

% ── regression cases for the adversarial-review fixes (S5–S10) ──────────────────

% (S6) MUST NOT FLAG — numeric-style citation [12] is a citation (dominant CS style).
Transformers outperform RNNs by 12\% [12].

% (S6b) MUST NOT FLAG — citation range/list [3, 4] and [1]--[3] both count.
Prior detectors reduce error by 20\% [3, 4].

% (S7) MUST NOT FLAG — biblatex \footcite is a citation macro.
Prior work improves accuracy by 15\% \footcite{smith2020}.

% (S7b) MUST NOT FLAG — \citeauthor / \parencite are citation macros too.
\citeauthor{jones2019} shows a 9\% reduction over the baseline \parencite{jones2019}.

% (S8) MUST NOT FLAG — "cf." must NOT sever the trailing \cite into its own segment.
RNN baselines underperform on long sequences, cf. \cite{vaswani2017}.

% (S8b) MUST NOT FLAG — "e.g." before a capital must not break the sentence off its cite.
Several detectors miss adversarial inputs, e.g. Gradient attacks \cite{goodfellow2015}.

% (S9) MUST NOT FLAG — float pointer via \ref-macro (no literal digit); verb refers to a Table.
Table~\ref{tab:results} shows the aggregated numbers across all runs.

% (S9b) MUST NOT FLAG — \autoref / \Cref float pointers, verb-only.
\Cref{fig:overview} demonstrates the pipeline end to end.

% (S10) MUST NOT FLAG — tabular body: numeric rows are data, not a prose claim.
The results appear below.
\begin{tabular}{lll}
Model & Acc & F1 \\
BERT  & 0.91 & 0.89 \\
GPT   & 0.94 & 0.93 \\
\end{tabular}

% (S5) MUST FLAG — bare country "US" must NOT be read as the paper's own result ("us").
US-based detectors reject 40\% of benign inputs.

% (S5b) MUST FLAG — "Our reading of the literature" is a claim about OTHERS, not an own result.
Our reading of the literature shows transformers outperform every prior baseline.

% (S5c) MUST NOT FLAG — genuine own result: "our method" is the subject.
Our method reduces false positives by 6\% over the strongest baseline.

\end{document}
